%=====================================================================
%  DarkTrace -- Experimental Methodology
%
%  Compiles standalone:  pdflatex -> bibtex -> pdflatex x2
%  Or drop the body into the main paper by copying everything between
%  \section{Experimental Methodology} and the bibliography, then delete
%  this file's preamble and \begin/\end{document}.
%=====================================================================
\documentclass[11pt,a4paper]{article}

\usepackage[utf8]{inputenc}
\usepackage[T1]{fontenc}
\usepackage{lmodern}
\usepackage{microtype}
\usepackage[margin=1in]{geometry}
\usepackage{amsmath,amssymb}
\usepackage{booktabs}
\usepackage{multirow}
\usepackage{array}
\usepackage{longtable}
\usepackage{graphicx}
\usepackage{xcolor}
\usepackage[hidelinks]{hyperref}
\usepackage{enumitem}
\usepackage{caption}

\newcommand{\framework}{\textsc{DarkTrace}}

\title{\framework{}: Experimental Methodology\\
\large Explainable and Forensically Verifiable Multilingual Dark Web\\
Threat Intelligence Framework}
\author{}
\date{}

\begin{document}
\maketitle

%=====================================================================
\section{Experimental Methodology}\label{sec:method}

\subsection{Design Philosophy and Overview}
The evaluation is organized around six research targets: (T1)~threat-text
classification, (T2)~multilingual detection, (T3)~severity-scoring accuracy,
(T4)~explainability quality, (T5)~evidence integrity (SHA-256 hashing and
chain-of-custody), and (T6)~comparison against baselines. Each target is treated as a
separate, pre-registered experiment with its own hypotheses, metrics, and statistical
tests, so that the integrated framework is justified component-by-component as well as
end-to-end. All data are drawn from \emph{sanitized, archived, public, or synthetic}
sources; no live, illegal, or unauthorized collection is performed
(Section~\ref{sec:method-ethics}). The pipeline, splits, seeds, and configurations are
released for reproducibility (Section~\ref{sec:method-repro}).

\paragraph{Hypotheses.}
\begin{description}[leftmargin=2.6em,style=nextline,itemsep=1pt]
  \item[H1] A fine-tuned multilingual transformer significantly outperforms classical
  ML and monolingual deep baselines on threat classification (T1, T6).
  \item[H2] Multilingual joint training narrows the cross-lingual gap; zero-shot
  transfer incurs a bounded, measurable penalty relative to monolingual upper bounds
  (T2).
  \item[H3] Hybrid severity scoring (classifier posterior $+$ interpretable features) is
  better calibrated and better ranked than classifier-confidence and static
  rule baselines (T3).
  \item[H4] SHAP and attention yield higher explanation fidelity than LIME on this text
  domain, and the three are complementary for analysts (T4).
  \item[H5] The integrity layer achieves $100\%$ tamper detection with $0\%$ false
  accepts at low throughput overhead (T5).
\end{description}

%---------------------------------------------------------------------
\subsection{Ethical Dataset Sources}\label{sec:method-ethics}
We use only datasets that are publicly released for research, archived, sanitized, or
synthetically generated. Table~\ref{tab:data} summarizes the layered strategy. Raw
illicit content is never re-published; only derived features, labels, and hashes are
released where licensing permits. Personally identifying information (PII) and victim
data are removed or pseudonymized before any processing.

\begin{table}[t]
\centering
\caption{Layered, ethical dataset strategy. ``Role'' indicates the target(s) each source
serves. Verify each source's licence and re-use terms before publication.}
\label{tab:data}
\small
\renewcommand{\arraystretch}{1.15}
\begin{tabular}{@{}p{3.4cm}p{5.7cm}p{2.6cm}@{}}
\toprule
\textbf{Layer / Source} & \textbf{Description} & \textbf{Role} \\
\midrule
Archived darknet corpora (e.g.\ sanitized Agora/forum archives)~\cite{sangher2023towards}
  & Publicly archived, de-identified marketplace listings and forum posts used in prior
    peer-reviewed work & T1, T3, T6 \\
CTI NER corpora (APTNER, DNRTI, CyNER, CyberNER)~\cite{wang2022aptner,evangelatos2021ner,echchammakhy2025cyberner,feng2025promptbart}
  & STIX-aligned, public annotated CTI text & actor/entity, T1 \\
Open-source CTI alert datasets (e.g.\ CySecAlert)~\cite{riebe2021cysecalert}
  & Public OSINT security-event text with relevance labels & T1, T6 \\
Public vulnerability text (NVD/CVE descriptions)
  & Standardized severity-bearing descriptions for severity supervision & T3 \\
Multilingual augmentation (machine-translated subsets $+$ native low-resource crime
  text)~\cite{alharbi2025bert,nazir2025multilingual}
  & EN, HI, AR, RU, FR/DE evaluation sets & T2 \\
Synthetic threat text (template- and LLM-generated, human-reviewed)
  & Controlled, fully shareable data for low-resource languages and rare classes & T2, T3 \\
\bottomrule
\end{tabular}
\end{table}

\paragraph{Synthetic data protocol.} To address low-resource languages and rare/harmful
classes without collecting illegal material, we generate synthetic threat-style text
from templates and instruction-tuned models, conditioned on \emph{abstracted} threat
patterns (never operational attack instructions). Every synthetic item is
human-reviewed, labeled, and flagged as synthetic so its effect can be ablated. Synthetic
data is used to \emph{augment}, never to replace, the real evaluation sets, and all
headline results are reported on real (non-synthetic) test data.

%---------------------------------------------------------------------
\subsection{Threat Class Labels}\label{sec:method-labels}
We adopt a flat, mutually informative taxonomy of threat categories plus a non-threat
class (Table~\ref{tab:labels}), chosen to be coverage-complete for dark-web CTI while
remaining annotatable with high agreement. Severity is modeled \emph{orthogonally} as an
ordinal band (Section~\ref{sec:method-severity}), so a single category (e.g.\ ``data
leak'') can span severities.

\begin{table}[t]
\centering
\caption{Threat-category labels (classification) and severity bands (scoring).}
\label{tab:labels}
\small
\renewcommand{\arraystretch}{1.15}
\begin{tabular}{@{}p{3.4cm}p{8.6cm}@{}}
\toprule
\textbf{Threat category} & \textbf{Definition (annotation guideline summary)} \\
\midrule
Malware / Ransomware      & Offering, sharing, or discussing malicious code, builders, or ransomware operations \\
Exploit / Vulnerability   & Sale or disclosure of exploits, 0-days, PoCs, CVE-linked weaknesses \\
Data / Credential Leak    & Trading or dumping stolen credentials, databases, PII \\
Fraud / Carding           & Payment-card fraud, money laundering, scam services \\
Hacking / Attack Services & Attack-as-a-service, botnets, DDoS-for-hire, access brokering \\
Other Illicit Trade       & Non-cyber illicit goods/services (retained as a coarse class; may be excluded) \\
Benign / Non-threat       & Off-topic, social, or non-malicious content (negative class) \\
\midrule
\textbf{Severity band} & \textbf{Criterion} \\
\midrule
Critical / High / Medium / Low & Ordinal scale from analyst rubric over impact, exploitability, target sensitivity, and recency/urgency \\
\bottomrule
\end{tabular}
\end{table}

\paragraph{Annotation.} Two or more trained annotators label each item for category and
severity following a written codebook; we report inter-annotator agreement (Cohen's
$\kappa$ / Krippendorff's $\alpha$) and adjudicate disagreements. Class imbalance is
expected and handled at training and evaluation time (macro-averaging,
class-balanced loss).

%---------------------------------------------------------------------
\subsection{Preprocessing}\label{sec:method-prep}
Preprocessing is identical across models to ensure fair comparison and is logged as
hashed, derived artifacts (T5). Steps:
\begin{enumerate}[leftmargin=1.5em,itemsep=1pt]
  \item \textbf{Normalization:} Unicode/encoding normalization, HTML/boilerplate
  stripping, whitespace and case handling (case preserved for transformers).
  \item \textbf{Noise handling:} removal of markup, navigation, and signature blocks;
  optional defanging of live IOCs (e.g.\ \texttt{hxxp}) without losing the token signal.
  \item \textbf{Deduplication:} exact and near-duplicate removal via MinHash/SimHash to
  prevent train/test leakage from reposted content.
  \item \textbf{Language identification:} per-document language tagging; routing to the
  native-multilingual path and, in parallel, a translate-then-classify path.
  \item \textbf{Translation (parallel path):} machine translation of non-English text to
  English, retained as a separate derived artifact so its effect is ablatable.
  \item \textbf{Tokenization/representation:} TF-IDF / $n$-grams for classical models;
  subword tokenization for transformers~\cite{devlin2019bert,conneau2020xlmr}.
  \item \textbf{Feature extraction (severity):} interpretable threat features
  (CVE/exploit mentions, credential counts, target-sensitivity cues, urgency terms,
  actor-reputation signals) computed for the hybrid scorer.
\end{enumerate}

%---------------------------------------------------------------------
\subsection{Models}\label{sec:method-models}
\paragraph{Baseline ML.} TF-IDF features with Logistic Regression, Linear SVM, Random
Forest, and XGBoost---fast, interpretable lower bounds and the T6 reference point.
\paragraph{Deep (non-transformer).} CNN, BiLSTM, and BiLSTM-CRF (the CRF variant for
sequence/NER tasks)~\cite{kim2020automatic}.
\paragraph{Monolingual transformers.} BERT/RoBERTa per language as per-language upper
bounds~\cite{devlin2019bert,sangher2023towards}.
\paragraph{Multilingual transformers.} mBERT and XLM-R (core
classifier)~\cite{conneau2020xlmr,vaswani2017attention}, evaluated under joint
multilingual, zero-shot, and few-shot transfer regimes~\cite{thangaraj2024crosslingual}.
\paragraph{Generative/LLM baselines.} An instruction-tuned LLM evaluated zero-/few-shot
for classification and extraction~\cite{shah2025madcti,sorokoletova2024scalable}.
\paragraph{Severity head (proposed).} A calibrated hybrid mapping
$\{$posterior, interpretable features$\}\rightarrow s\in[0,1]$ and an ordinal band,
compared against (i)~raw classifier confidence and (ii)~a static rule/CVSS-style
baseline~\cite{zeng2022licality,saklani2025severity}.

%---------------------------------------------------------------------
\subsection{Train/Test Splits and Protocol}\label{sec:method-splits}
\begin{itemize}[leftmargin=1.5em,itemsep=1pt]
  \item \textbf{Primary split:} stratified 70/15/15 train/validation/test, preserving
  class proportions, with fixed seeds.
  \item \textbf{Cross-validation:} stratified $5$-fold CV for classical and deep models;
  for transformers, $3$ seeds $\times$ the held-out test set, reporting mean $\pm$ std.
  \item \textbf{Cross-lingual protocol (T2):} (a)~\emph{monolingual} (train/test same
  language) as upper bound; (b)~\emph{zero-shot} (train EN, test HI/AR/RU/FR/DE);
  (c)~\emph{few-shot} ($k$ target-language examples added). We report the transfer gap
  $\Delta\text{F1}$.
  \item \textbf{Leakage control:} deduplication before splitting; near-duplicates and
  reposts forced into the same split; language-disjoint and (where timestamps exist)
  \emph{temporal} splits (train on older, test on newer) to test robustness to drift.
  \item \textbf{Synthetic isolation:} synthetic data may enter training only; all
  headline test sets are real.
\end{itemize}

%---------------------------------------------------------------------
\subsection{Severity Scoring Evaluation}\label{sec:method-severity}
Severity is evaluated as both \emph{ranking} and \emph{calibration}: ranking via
NDCG@$k$ and Spearman/Kendall correlation against the analyst-assigned ordering;
calibration via Expected Calibration Error (ECE), Brier score, and reliability diagrams,
with temperature scaling applied post-hoc. Ordinal-band performance uses
macro-F1 and Quadratic Weighted Kappa (which penalizes distant misclassifications more).

%---------------------------------------------------------------------
\subsection{Explainability Evaluation}\label{sec:method-xai}
We compare SHAP~\cite{lundberg2017shap}, LIME~\cite{ribeiro2016lime}, and transformer
attention, following comparative-XAI practice in
security~\cite{hermosilla2025xai,gaspar2024xai,kalakoti2025evaluating,salih2023perspective}:
\begin{itemize}[leftmargin=1.5em,itemsep=1pt]
  \item \textbf{Fidelity:} deletion/insertion and AOPC (area over the perturbation
  curve)---removing high-attribution tokens should drop confidence fastest for the most
  faithful method.
  \item \textbf{Stability/robustness:} attribution agreement under small input
  perturbations and across seeds.
  \item \textbf{Complexity/sparsity:} number of features needed to explain a decision.
  \item \textbf{Cross-method agreement:} SHAP--LIME--attention rank correlation.
  \item \textbf{Cross-lingual fidelity:} whether explanation quality degrades under
  translation/transfer (links T2 and T4).
\end{itemize}

%---------------------------------------------------------------------
\subsection{Evidence Integrity Evaluation}\label{sec:method-integrity}
The integrity layer is validated with an adversarial \emph{tamper test}. A controlled
set of post-hoc modifications---single/multi bit-flips, token substitution, artifact
deletion, custody-record reordering, and ledger edits---is injected, and verification
recomputes $\text{SHA-256}$ digests and re-validates the hash-chained custody
log~\cite{lone2019forensicchain,ali2022procedure,loffi2025management}. Metrics:
tamper-detection rate (target $100\%$), false-accept rate (target $0\%$), verification
latency, custody-log completeness, and the throughput overhead the layer imposes on the
analytic pipeline.

%---------------------------------------------------------------------
\subsection{Evaluation Metrics Summary}\label{sec:method-metrics}
\begin{table}[t]
\centering
\caption{Metrics by research target.}
\label{tab:method-metrics}
\small
\renewcommand{\arraystretch}{1.15}
\begin{tabular}{@{}p{3.0cm}p{9.0cm}@{}}
\toprule
\textbf{Target} & \textbf{Metrics} \\
\midrule
T1 Classification & Accuracy, macro-/weighted-F1, per-class precision/recall/F1, MCC \\
T2 Multilingual   & Per-language F1, macro-F1 across languages, transfer gap $\Delta$F1, language-wise confusion \\
T3 Severity       & NDCG@$k$, Spearman/Kendall $\rho$, ECE, Brier, Quadratic Weighted $\kappa$ \\
T4 Explainability & Fidelity (AOPC), stability, complexity, cross-method agreement \\
T5 Integrity      & Tamper-detection rate, false-accept rate, verification latency, log completeness, overhead \% \\
T6 Baselines      & All of the above, head-to-head, with statistical tests (Sec.~\ref{sec:method-stats}) \\
\bottomrule
\end{tabular}
\end{table}

%---------------------------------------------------------------------
\subsection{Ablation Study}\label{sec:method-ablation}
We remove or alter one capability at a time and measure the effect on the most-affected
axis (Table~\ref{tab:method-ablation}), establishing that \emph{integration} is the
contribution rather than any single module.

\begin{table}[t]
\centering
\caption{Ablation configurations (result cells filled from experiments).}
\label{tab:method-ablation}
\small
\renewcommand{\arraystretch}{1.15}
\begin{tabular}{@{}p{6.0cm}p{4.2cm}c@{}}
\toprule
\textbf{Configuration} & \textbf{Primary axis} & \textbf{$\Delta$ vs.\ full} \\
\midrule
Full \framework{}                              & ---                       & --- \\
$-$ multilingual (English-only)                & T2 cross-lingual F1       & \emph{$-$} \\
\;native vs.\ translate-then-classify          & T2 cross-lingual F1       & \emph{$-$} \\
$-$ hybrid features (posterior only)           & T3 calibration (ECE)      & \emph{$-$} \\
$-$ calibration (no temperature scaling)       & T3 ECE/Brier              & \emph{$-$} \\
SHAP only / LIME only / attention only         & T4 fidelity               & \emph{$-$} \\
$-$ synthetic augmentation                     & T1/T2 low-resource F1     & \emph{$-$} \\
$-$ integrity layer                            & T5 tamper detection       & \emph{$-$} \\
$-$ STIX extraction / link analysis            & actor recall / usability  & \emph{$-$} \\
\bottomrule
\end{tabular}
\end{table}

%---------------------------------------------------------------------
\subsection{Tables and Graphs to Include}\label{sec:method-figs}
\textbf{Tables:} (1)~dataset summary with per-language/per-class counts; (2)~main
classification results (Acc/F1/MCC) across all models; (3)~per-language and transfer-gap
table; (4)~severity ranking + calibration table; (5)~explanation-quality table
(fidelity/stability/agreement); (6)~integrity tamper-test table; (7)~ablation table;
(8)~statistical-test summary (p-values, effect sizes).
\textbf{Figures:} (a)~architecture/data-flow diagram; (b)~per-class confusion matrices
(overall and per language); (c)~reliability diagram for severity calibration;
(d)~AOPC/deletion curves comparing SHAP/LIME/attention; (e)~example SHAP/attention
token-attribution heatmaps (EN and one non-EN language); (f)~critical-difference (CD)
diagram from the Friedman--Nemenyi test; (g)~actor link-analysis graph;
(h)~integrity-overhead vs.\ throughput plot.

%---------------------------------------------------------------------
\subsection{Statistical Validation}\label{sec:method-stats}
All comparisons report central tendency with dispersion and significance:
\begin{itemize}[leftmargin=1.5em,itemsep=1pt]
  \item \textbf{Pairwise classifier comparison:} McNemar's test on the shared test
  set~\cite{mcnemar1947note}; the $5\times2$~cv paired $t$-test for
  resampled comparisons~\cite{dietterich1998approximate}.
  \item \textbf{Multiple models over multiple datasets/languages:} Friedman test with
  Nemenyi post-hoc and a critical-difference diagram~\cite{demsar2006statistical}.
  \item \textbf{Confidence intervals:} bootstrap (e.g.\ $1000$ resamples) 95\% CIs for
  F1, NDCG, and ECE.
  \item \textbf{Effect sizes and correction:} report effect sizes (e.g.\ Cliff's
  $\delta$) and apply Holm/Bonferroni correction for multiple comparisons.
  \item \textbf{Calibration significance:} compare ECE/Brier with bootstrap CIs across
  scoring variants.
  \item Significance threshold $\alpha = 0.05$ unless corrected.
\end{itemize}

%---------------------------------------------------------------------
\subsection{Reproducibility Checklist}\label{sec:method-repro}
\begin{itemize}[leftmargin=1.5em,itemsep=1pt]
  \item[$\square$] Public release of code, configs, and trained-model hashes.
  \item[$\square$] Dataset datasheet; construction scripts and exact split files;
  SHA-256 of each released artifact.
  \item[$\square$] Fixed random seeds; reported hardware, library versions, and
  environment file.
  \item[$\square$] Hyperparameter search space and selected values; training/inference
  cost (time, energy).
  \item[$\square$] Exact preprocessing, tokenization, and translation settings.
  \item[$\square$] Per-run metrics (not just best run); mean~$\pm$~std over seeds/folds.
  \item[$\square$] Statistical-test scripts and raw outputs.
  \item[$\square$] Ethics statement, data-management plan, and licence compliance for
  every source.
  \item[$\square$] Synthetic-data generation prompts/templates and human-review records.
\end{itemize}

%---------------------------------------------------------------------
\subsection{Limitations}\label{sec:method-limits}
(1)~Archived/sanitized and synthetic data may underrepresent the most recent, ephemeral
threats; live coverage is constrained by ethics and law. (2)~Severity labels are partly
subjective; inter-annotator agreement bounds calibration claims. (3)~Machine translation
can distort low-resource semantics and the explanations derived from translated text---an
effect we measure but cannot fully remove. (4)~Synthetic augmentation risks distributional
mismatch; we mitigate by testing only on real data and ablating synthetic contributions.
(5)~The permissioned-ledger integrity model resists post-hoc tampering but assumes
trusted timestamping and does not prevent collection-time fabrication. (6)~Adaptive
adversaries (obfuscation, code-switching, slang drift) may evade detection, motivating the
temporal-split robustness analysis.

%=====================================================================
\bibliographystyle{plain}
\bibliography{references}

\end{document}
